\documentclass[11pt,a4paper]{article}
\usepackage{agpstyle}

\title{\bfseries A Three-Stage PPO Curriculum for the\\ Continuous Art Gallery Problem\\[4pt]
\large Document 2 of 2: \stageI{} $\rightarrow$ \stageII{} $\rightarrow$ \stageIII{}}
\author{AGP Project --- Design Note}
\date{\today}

\begin{document}
\maketitle

\begin{abstract}
\noindent
Document 1 defined a single MDP for the Art Gallery Problem in which coverage and guard
parsimony enter as two potential-based reward components. This document specifies the
training curriculum on top of it. Stage~\stageI{} (\stageI{}) switches the guard cost off
entirely and asks only for coverage, with a boundary-inward placement prior that is
annealed away; the policy learns to fill the gallery, deliberately over-completely.
Stage~\stageII{} (\stageII{}) unmasks a delete action, switches the guard cost on, and pays
a bounded bonus for every removal that costs no coverage, decaying smoothly as coverage is
lost. Stage~\stageIII{} (\stageIII{}) replaces the fixed regularity penalties with
state-wise constraints and polishes on the hardest instances. The design borrows four
specific devices from the humanoid locomotion literature --- one reward basis with a
per-stage weight column, a task objective held at exactly zero weight until its
prerequisite is learned, scaffolding terms deleted rather than decayed, and multi-critic
advantage weighting --- and contributes three things those papers do not provide: a
quantitative stage-transition gate, a blended switch with a shadow-trained parsimony critic,
and a cheap privileged \emph{prune oracle} that can warm-start the delete action by
imitation. We give the full weight schedule, PPO hyperparameters per stage, network
architecture, the instance-difficulty curriculum, an ablation protocol, and a concrete
implementation roadmap against the existing code.
\end{abstract}

\tableofcontents

% ===========================================================================
\section{Why stage the training at all}
% ===========================================================================

\subsection{The conflict}

The AGP objective has two parts that pull against each other during \emph{learning}, though
not at the optimum. Coverage is easy to improve from a random start --- almost any guard
placed in an empty gallery captures a large fraction of the area --- and its gradient is
dense and informative. Parsimony is the opposite: a guard-count penalty is meaningful only
once the agent can reliably cover the gallery, because before that point the cheapest way
to reduce the guard count is to stop placing guards, which also destroys coverage. Turning
both on from step zero means the agent must discover, from a random policy, that the
\emph{only} route to a high return runs through a temporarily expensive regime.

This is the same structure \citet{wang2025dribble} describe for humanoid dribbling:

\begin{quote}\small
``Training a humanoid robot to dribble presents a fundamental challenge: the rewards
associated with walking and dribbling often conflict and are difficult to balance. Walking
requires the robot to maintain stability and move efficiently, while dribbling demands
agility and precise ball control. This trade-off complicates the simultaneous optimization
of both skills.''
\end{quote}

\noindent
Their ablation quantifies the cost of ignoring it: a single-stage policy reaches $33.0\%$
success on dribble-to-target and $0\%$ on obstacle avoidance, against $86.7\%$ and $92.3\%$
for the two-stage curriculum. The failure mode they report --- ``one-stage training policy
can't stably walk and control the ball simultaneously, as both capabilities are learned
jointly from scratch'' --- is exactly the failure mode we expect from an AGP policy trained
with the guard cost on from the beginning.

\subsection{What the three reference papers actually do}

\Cref{tab:papers} summarises the three papers, with an eye to what is reusable.

\begin{table}[t]
\centering
\small
\caption{The three locomotion references, compared on the dimensions that matter here.
Local copies are in \code{references/locomotion/}.}
\label{tab:papers}
\begin{tabularx}{\linewidth}{@{}p{2.5cm} p{2.1cm} X p{2.5cm}@{}}
\toprule
\textbf{Paper} & \textbf{Stages} & \textbf{What changes at a stage boundary} & \textbf{Transition trigger} \\
\midrule
Dribble Master \cite{wang2025dribble}
& 2
& Task reward $0.0 \to 1.5$; approach proxy $2.0 \to 1.0$; imitation scaffold $1.2 \to 0.0$;
  ball spawn $10$\,m $\to 2$\,m; camera FOV $2\times \to 1\times$ real
& Qualitative only: ``once the robot consistently demonstrates reliable locomotion'' \\
\addlinespace
Vision-driven soccer \cite{wang2025vision}
& 1 (explicitly single-stage)
& --- (no curriculum; uses potential-based shaping and a multi-critic split instead)
& n/a \\
\addlinespace
Agile striker \cite{xu2025striker}
& 4
& Task reward $0.0 \to 4.0$; tracking proxy $2.0 \to 0.5$; gait scaffolds $3.0/1.0 \to 0.0$;
  conflicting regulariser (feet yaw) $-1.0 \to 0.0$; spawn annulus $[0.5,4]\to[0.5,2]$\,m;
  then DAgger distillation; then constrained online RL
& Qualitative only: ``once the teacher policy acquires a basic kicking skill'' \\
\bottomrule
\end{tabularx}
\end{table}

Four patterns are worth importing wholesale.

\begin{enumerate}[label=\textbf{P\arabic*.},leftmargin=2.7em]
  \item \textbf{One reward basis, one weight column per stage.} Neither multi-stage paper
        changes the observation space, the action space or the reward \emph{terms} between
        stages --- only the coefficients (and the initial-state distribution). We adopt this
        exactly: \cref{tab:schedule} is the entire curriculum.
  \item \textbf{The objective you cannot yet optimise gets weight zero, not a small weight.}
        Dribble Master holds ball-velocity tracking at $0.0$ through stage 1;
        the striker paper holds its ball-velocity reward at $0.0$ through stage 1. Neither
        uses a small-but-nonzero coefficient. Our guard cost is therefore \emph{exactly}
        $0$ in \stageI{} --- which is also what was asked for in the project brief.
  \item \textbf{Scaffolds are deleted; proxies are demoted.} Both papers zero their
        imitation/gait scaffolding at the boundary (Dribble Master: reference joint
        position $1.2 \to 0.0$, ``removed to promote exploration and discovery of effective
        dribbling behaviors''; striker: feet swing $3.0 \to 0.0$, arm swing $1.0 \to 0.0$),
        but keep the stage-1 proxy alive at a reduced weight ($2.0 \to 1.0$ and
        $2.0 \to 0.5$ respectively) rather than removing it, as insurance against
        catastrophic forgetting. Document~1's \emph{anneal-iff-non-potential} rule explains \emph{why} this
        split is the right one: scaffolds are non-potential and bias the optimum, proxies
        closer to the objective do not.
  \item \textbf{Split the critic before you split the reward.}
        \citet{wang2025vision} found that ``using a single critic to estimate combined
        rewards can cause negative interference between distinct reward components,'' and
        use $\Adv_{\mathrm{total}} = 2\Adv_{\mathrm{goal}} + 1\cdot\Adv_{\mathrm{aux}}$. Our
        coverage and parsimony rewards differ by more than an order of magnitude in scale
        and by a factor of $\sim\!30$ in per-step variance, so this applies with force.
\end{enumerate}

\subsection{What we have to design ourselves}

None of the three papers gives a quantitative criterion for advancing a stage. Dribble
Master says ``once the robot consistently demonstrates reliable locomotion and target
pursuit''; the striker paper says ``once the teacher policy acquires a basic kicking
skill''. There is no threshold, no step count, and no automated gate anywhere. Nor does any
of them address the value-function discontinuity that a hard reward switch creates.
\Cref{sec:transitions} supplies both. A third gap --- how to get a policy to \emph{use} a
newly unmasked action that it has never sampled --- is addressed in \cref{sec:oracle} with a
privileged prune oracle, which is the AGP analogue of the teacher--student step in
\citet{xu2025striker} but far cheaper, because unlike a robot teacher our oracle is a
twenty-line greedy loop.

% ===========================================================================
\section{Curriculum overview}
% ===========================================================================

\begin{figure}[h]
\centering
\begin{tikzpicture}[font=\small, x=1cm, y=1cm]
  % --- gate markers (drawn first so boxes sit on top)
  \foreach \x/\g in {4.9/1, 9.6/2}{
    \draw[agpgreen, thick, dashed] (\x,-0.35) -- (\x,3.15);
    \node[font=\scriptsize\bfseries, agpgreen, above] at (\x,3.15) {gate $G_{\g}$};
  }

  % --- stage boxes
  \node[draw, thick, fill=agpgrey, rounded corners, minimum height=2.1cm, text width=4.1cm,
        align=center, anchor=south west] at (0.25,0.55)
    {\textbf{\stageI{}}\\[3pt]
     \scriptsize $\ADD$ only\\
     \scriptsize $w_{\mathrm{par}} = 0$ (no guard cost)\\
     \scriptsize sweep $+$ frontier priors, annealed\\[2pt]
     \scriptsize\itshape\color{agpaccent} ends with I-b: over-saturate};
  \node[draw, thick, fill=agpgrey, rounded corners, minimum height=2.1cm, text width=4.1cm,
        align=center, anchor=south west] at (5.05,0.55)
    {\textbf{\stageII{}}\\[3pt]
     \scriptsize $\ADD$ / $\DEL$ / $\STOP$\\
     \scriptsize $w_{\mathrm{par}} = 0.02$, prune bonus on\\
     \scriptsize scaffolds at zero\\[2pt]
     \scriptsize\itshape\color{agpaccent} prune-oracle warm start};
  \node[draw, thick, fill=agpgrey, rounded corners, minimum height=2.1cm, text width=3.7cm,
        align=center, anchor=south west] at (9.75,0.55)
    {\textbf{\stageIII{}}\\[3pt]
     \scriptsize state-wise constraints\\
     \scriptsize $+\,\MOVE$ action\\
     \scriptsize hard / extreme tiers};

  % --- timeline axis
  \draw[-{Latex[length=3mm]}, thick] (0,0) -- (14.4,0);
  \node[right, font=\scriptsize] at (14.4,0) {steps};
  \foreach \x/\lab in {0.25/0, 4.9/1.5M, 9.6/3.0M, 13.7/3.75M}{
    \draw (\x,0.14) -- (\x,-0.14);
    \node[below, font=\scriptsize] at (\x,-0.16) {\lab};
  }
\end{tikzpicture}
\caption{The curriculum. Step counts are targets, not fixed schedules: each boundary is
crossed when its performance gate (\cref{sec:transitions}) fires, with the step count acting
only as a fallback cap. Weights are linearly interpolated across a $50$k-step blend window
rather than switched instantaneously.}
\label{fig:timeline}
\end{figure}

% ===========================================================================
\section{Stage I --- \stageI{}}
% ===========================================================================

\subsection{Objective and configuration}

\stageI{} answers one question: \emph{can the policy cover an arbitrary gallery at all?}
The guard cost is zero, so the telescoped episode objective (Document 1,
the \emph{telescoping proposition}) reduces to
\[
  R_{\stageI{}}(\tau) \;=\; w_{\mathrm{cov}}\,c_T \;+\; w_{\mathrm{goal}}\Ind{c_T \ge c^{\dagger}}
  \;-\; w_{\mathrm{fail}}\Ind{c_T < c^{\dagger}}
  \;+\; \text{(scaffolding)} ,
\]
i.e.\ pure maximum coverage. This is a problem with a known, strong, cheap baseline: greedy
maximum coverage, which inherits the $1 - 1/e$ guarantee from submodularity
\cite{nemhauser1978}. \stageI{} is not expected to beat greedy. It is expected to
\emph{match} greedy while (a) using continuous positions rather than a candidate set and
(b) producing, as a side effect, exactly the over-complete guard sets that \stageII{} needs.

\begin{description}[style=nextline,leftmargin=1.2em]
\item[Action mask:] mode restricted to $\{\ADD\}$. This makes the action space identical to
  the current $[0,1]^2$ implementation.
\item[Instance distribution $\Dctx$:] all four difficulty tiers, but sampled with a
  difficulty-annealed weighting (\cref{sec:instances}): $(0.55, 0.30, 0.10, 0.05)$ over
  (easy, mid, hard, extreme) at step $0$, linearly reaching uniform by the end of the stage.
\item[Termination:] at $c_t \ge c^{\dagger}$, on the guard budget $K_{\max} = 24$, or on the
  patience stall.
\item[Discount:] $\gamma = 0.99$.
\end{description}

\subsection{The boundary-inward sweep prior}

The project brief proposes an explicit placement order: ``go around the edges and come
inwards''. This is a good prior for the AGP, because occlusion in a polygon with holes is
generated at the boundary --- specifically at reflex vertices, which are what make
Chv\'atal's $\lfloor n/3 \rfloor$ bound tight \cite{chvatal1975,fisk1978,orourke1987} --- and a
guard placed against a wall sees a half-plane-like region with no wasted coverage behind it,
whereas an interior guard's visibility polygon is partly redundant with itself.

We encode it as the term $r_{\mathrm{sweep}}$ of Document 1, a bounded exponential kernel on
the distance-to-boundary $\dbd$ whose target radius sweeps inward \emph{within} an episode:
\[
  r_{\mathrm{sweep}}(\xi, t) \;=\;
  \exp\!\left(-\frac{\bigl(\dbd(\xi) - \mu_{\mathrm{ring}}(t)\bigr)^2}{2\sigma_s^2}\right),
  \qquad
  \mu_{\mathrm{ring}}(t) = \Bigl(\mu_0 + (\mu_1-\mu_0)\tfrac{t}{K_{\max}}\Bigr) \cdot L ,
\]
with $L$ the bounding-box diagonal, $\mu_0 = 0.02$, $\mu_1 = 0.25$, $\sigma_s = 0.05 L$.
At $t=0$ the policy is paid for placing within $2\%$ of the diagonal of a wall; by the end
of the guard budget the preferred band has moved to $25\%$, i.e.\ well into the interior.

Two properties matter. First, the kernel is bounded in $(0,1]$, so its weight \emph{is} its
maximum per-step contribution --- the convention every tracking reward in
\citet{wang2025dribble} and \citet{xu2025striker} follows, and the reason a dozen terms can be
balanced by hand. Second, it is a prior on the placement \emph{order}, not on the final
configuration, which is structurally the same thing as the gait-phase indicators
$\mathbf{1}_{\mathrm{feet\ swing}} \cdot \mathbf{1}_{\mathrm{swing\ period}}$ that the striker
paper uses to install a walking rhythm in stage 1 and then zeroes in stage 2.

\takeaway{%
$r_{\mathrm{sweep}}$ is \emph{not} potential-based and therefore biases the optimum
(Document~1, \emph{anneal-iff-non-potential}): a wall-hugging guard set is not in general optimal.
Its weight must reach $0$ before \stageI{} ends. We ramp it down linearly over the final
$30\%$ of the stage, so that the last part of \stageI{} is trained under the unbiased
objective and the policy has time to unlearn the prior's distortions before the gate is
evaluated.}

The same applies to $r_{\mathrm{front}}$ (pay for landing in still-uncovered space) and
$r_{\mathrm{sep}}$ (punish placing within $d_{\min}$ of an existing guard). Both are useful
early --- they are what stop a random policy from stacking guards on one spot --- and both
are wrong late: in a narrow corridor two guards genuinely should sit close together, and a
guard placed in already-covered space can still be the best next move if it opens a new
region. They anneal on the same schedule.

\subsection{Stage I-b: deliberate over-saturation}\label{sec:overs}

For the final $\sim\!150$k steps of \stageI{} we change one thing: the episode no longer
terminates at $c^\dagger$, and instead runs to
$K_{\mathrm{sat}} = \min\bigl(K_{\max}, \lceil 1.4\, k_{\mathrm{ref}} \rceil\bigr)$ guards.
The purpose is to manufacture the state distribution that \stageII{} will be trained on.
Without this, the policy hands \stageII{} a guard set that is already close to minimal,
almost no deletion is free, and the delete head learns nothing.

This is the AGP version of a trick both reference papers use in different guises: make the
signal you want the agent to learn from artificially abundant early.
\citet{wang2025dribble} widen the simulated camera FOV to \emph{twice} the real
RealSense D455 field of view during stage 1, so that the active-sensing reward fires densely,
then shrink it to the true value in stage 2. We inflate the guard count so that the
redundancy signal fires densely, then let \stageII{} deflate it.

By the \emph{redundancy-monotonicity proposition} of Document~1, the redundancy deficits $\delta^-_j$ measured at
$K_{\mathrm{sat}}$ guards are the smallest they will ever be, so this is also the state in
which the pruning decision is least ambiguous.

\subsection{Exit gate $G_1$}

See \cref{sec:transitions}. Informally: held-out feasibility $\ge 0.90$ on easy/mid and
$\ge 0.75$ on hard/extreme, sustained over three consecutive evaluation rounds, with all
scaffolding weights already at zero.

% ===========================================================================
\section{Stage II --- \stageII{}}
% ===========================================================================

\subsection{What changes}

Three things, all of them in the weight column and the action mask:

\begin{enumerate}[label=(\alph*)]
  \item the mode head is unmasked to $\{\ADD, \DEL, \STOP\}$;
  \item the guard-count potential is switched on, $w_{\mathrm{par}}: 0 \to 0.02$, which by
        construction (Document 1) means $-0.02$ per $\ADD$ and $+0.02$ per $\DEL$ --- exactly
        the ``adding is negative, removing is positive'' structure asked for;
  \item the clean-prune bonus $w_{\mathrm{prune}}\exp(-\delta^{-}/\sigma_p)$ is switched on,
        which pays close to the full $0.15$ when a deletion costs no coverage and decays to
        nothing as coverage is lost --- the ``lower reward scaling with space lost'' clause.
\end{enumerate}

Simultaneously all three scaffolding weights are already at zero, and the terminal
parsimony bonus $w_{\mathrm{pars}}(k_{\mathrm{ref}} - k_T)$ turns on at $0.25$.
The coverage potential does \emph{not} change: it is potential-based, so leaving it at
$w_{\mathrm{cov}} = 1.0$ costs nothing in bias and buys full protection against forgetting
how to cover. This is where our design departs from \citet{wang2025dribble} and
\citet{xu2025striker}, both of whom must demote their stage-1 proxy ($2.0 \to 1.0$ and
$2.0 \to 0.5$) precisely because theirs is \emph{not} potential-based and would otherwise
compete with the real objective.

\subsection{The delete action will not be discovered by accident}\label{sec:oracle}

This is the main practical risk of the whole curriculum, and it has no counterpart in the
reference papers, whose action space never changes. At the moment $\DEL$ is unmasked, the
policy has taken millions of $\ADD$ actions and zero $\DEL$ actions. Its mode-head logits
for $\DEL$ and $\STOP$ have received no meaningful gradient. Worse, the first deletions it
tries will be uniformly distributed over guards and will therefore usually be expensive,
producing negative advantage and pushing the $\DEL$ logit further down. The likely outcome
of a naive switch is a policy that never deletes and \stageII{} that does nothing.

Four mitigations, in increasing order of effort.

\paragraph{1. Reset the mode head.} Re-initialise the three mode logits to zero (uniform)
at the switch, keeping the shared trunk and the point head. One line, and it removes the
accumulated bias against unsampled modes.

\paragraph{2. A mode-specific entropy floor.} Apply the entropy bonus to the categorical
mode distribution separately from the continuous point distribution, with
$\beta_{\mathcal{H}}^{\mathrm{mode}} = 0.01$ decaying to $0.002$ over the first $300$k steps
of \stageII{} and $\beta_{\mathcal{H}}^{\mathrm{point}} = 0.003$ throughout. The current code
runs \code{ENT\_COEF = 0.0}, which is fine for a single continuous head but wrong for a
discrete head that must explore a newly available option.

\paragraph{3. A decaying logit bias, not forced actions.} It is tempting to force
$\DEL$ with some probability $\varepsilon$ during early \stageII{}. Do not: PPO's importance
ratio $\pi_\theta(a|s)/\pi_{\theta_{\mathrm{old}}}(a|s)$ is computed against the policy that
is assumed to have generated the data, and injecting off-policy actions silently corrupts
it. Instead add a bias $b(t)$ to the $\DEL$ logit \emph{inside} the policy, so the behaviour
policy and the policy being scored are the same distribution, and anneal $b$ from $1.5$ to
$0$ over $200$k steps.

\paragraph{4. A privileged prune oracle with DAgger warm-start.} The strongest option, and
cheap here in a way it is not in robotics. Define:

\begin{algorithm}[h]
\small
\caption{\textsc{GreedyPruneOracle}($\Gallery$, $\Guards$, $c^\dagger$)}
\begin{algorithmic}[1]
\While{true}
  \State $\delta^{-}_j \gets f(\Guards) - f(\Guards \ominus g_j)$ \textbf{ for all } $g_j \in \Guards$
  \State $j^\star \gets \arg\min_j \delta^{-}_j$
  \If{$f(\Guards) - \delta^{-}_{j^\star} < c^{\dagger}$} \State \textbf{break} \EndIf
  \State $\Guards \gets \Guards \ominus g_{j^\star}$; \; \textbf{emit} $(\DEL, g_{j^\star})$
\EndWhile
\State \textbf{emit} $(\STOP, \cdot)$
\end{algorithmic}
\end{algorithm}

\noindent
One oracle call costs $O(k^2)$ Shapely unions, which for $k \le 24$ is milliseconds --- the
per-guard visibility polygons are already cached in \code{Security.all\_guards}. The oracle
is \emph{privileged} in the same sense as the teacher policies in
\citet{xu2025striker} and \citet{pinto2018asymmetric}: it uses exact continuous geometry that
the rasterised student observation does not expose.

Use it exactly as they use their teacher, with DAgger \cite{ross2011dagger}: roll out the
student, query the oracle for the correct action at each visited state, and add a
behaviour-cloning loss on the mode head and the point head,
\[
  \mathcal{L}_{\mathrm{BC}}
  \;=\; \Ex_{s \sim \pi_\theta}\Bigl[\,
    \mathrm{CE}\bigl(\mu_\theta(s),\, \mu^{\mathrm{oracle}}(s)\bigr)
    + \lVert \xi_\theta(s) - \xi^{\mathrm{oracle}}(s)\rVert^2 \,\Bigr],
\]
with coefficient $1.0$ for the first $200$k steps of \stageII{} and linearly to $0$ by
$500$k. Rolling out the \emph{student} and labelling with the oracle (rather than cloning
oracle trajectories) is the whole point of DAgger: the student is corrected on the states
its own mistakes produce. \citet{xu2025striker} state the same rationale --- ``the student
learns not only from the teacher's demonstrations but also from states encountered due to
its own errors''.

\takeaway{%
The oracle is a \emph{warm start}, not the target. Greedy pruning is myopic: it deletes the
cheapest guard at each step and never considers deleting two guards and re-adding one
better-placed guard. The policy has both $\DEL$ and $\ADD$ available and can learn exactly
that move, which is the main source of expected improvement over greedy. This is why
$\mathcal{L}_{\mathrm{BC}}$ must decay to zero --- the same reason
\citet{wang2025dribble} delete their imitation reward entirely rather than decay it, ``to
promote exploration and discovery of effective dribbling behaviors''.}

\subsection{Preventing coverage regression}

The coverage ratchet of Document 1, $c_t \ge \min(\bar c_t, c^\dagger) - \delta_{\mathrm{slack}}$,
is enforced in \stageII{} as an action mask: a $\DEL$ that would violate it is rejected and
charged $r_{\mathrm{void}}$. Note the ratchet is measured against the \emph{running maximum}
$\bar c_t$, not against the previous step, so the $1\%$ slack cannot be spent repeatedly ---
there is no drift exploit where the agent deletes a guard costing $0.9\%$ twenty times over.

\subsection{Episode initialisation}

\stageII{} episodes are drawn from a mixture:

\begin{center}
\small
\begin{tabularx}{\linewidth}{@{}l l X@{}}
\toprule
\textbf{Share} & \textbf{Initial state} & \textbf{Purpose} \\
\midrule
$60\%$ & empty gallery & the policy must still be able to solve an instance end to end \\
$40\%$ & saturated: $K_{\mathrm{sat}}$ guards from the frozen \stageI{} policy & dense pruning signal from step $0$ \\
\bottomrule
\end{tabularx}
\end{center}

The $40\%$ saturated share is the direct analogue of moving the ball spawn from $10$\,m to
$2$\,m \cite{wang2025dribble} or tightening the robot spawn annulus from $[0.5,4.0]$ to
$[0.5,2.0]$\,m \cite{xu2025striker}: start the agent where the skill being learned is
actually exercised. The $60\%$ empty share is the insurance those papers buy differently,
by demoting rather than deleting their stage-1 proxy.

\subsection{Exit gate $G_2$}

Feasibility not degraded by more than $2$ percentage points relative to the end of
\stageI{}, mean pruning yield $\ge 2$ guards, and competitive ratio
$\Ex[k_T/k_{\mathrm{ref}}] \le 1.0$ --- i.e.\ the policy is at least matching greedy on guard
count while holding coverage.

% ===========================================================================
\section{Stage III --- \stageIII{} (optional)}
% ===========================================================================

\stageIII{} is worth running only if \stageII{} converges and the residual failures are
concentrated on the \code{hard}/\code{extreme} tiers, which is the expected outcome. It
makes three changes.

\subsection{From penalties to state-wise constraints}

The regularity terms ($r_{\mathrm{proj}}$, $r_{\mathrm{void}}$, $r_{\mathrm{stop}}$) move out of
the reward and into a constraint,
\[
  \max_\pi \; \Ex\Bigl[\textstyle\sum_t \gamma^t r^{\mathrm{task}}_t\Bigr]
  \quad\text{s.t.}\quad
  c^{\mathrm{cost}}(s_t,a_t) \le h \quad \forall (s_t,a_t)\in\tau ,
\]
solved with a constrained policy-gradient method such as N-P3O \cite{lee2023constrained} or
CPO \cite{achiam2017cpo}; see also \citet{altman1999}. The motivation transfers exactly from
\citet{xu2025striker}, who diagnose \emph{inhomogeneous credit assignment}: their task reward
is concentrated at the kick, so ``a fixed regularization coefficient can impose
disproportionately large penalties at certain steps --- particularly those immediately
preceding the kick, where the agent expects high immediate future reward.''

AGP has the identical profile. The first guard in an empty gallery typically captures
$30$--$50\%$ of the area; the last one before $c^\dagger$ captures a sliver. A fixed
$-0.25$ aiming penalty is negligible against the first and dominant against the last, so it
suppresses exactly the delicate end-game placements that determine whether the instance is
solved at all. Their ablation is the reason to take this seriously: raising the fixed
coefficient from $1.0$ to $1.5$ \emph{dropped} success from $64.8\%$ to $43.4\%$, while the
constrained formulation reached $79.5\%$. That is the signature of a problem scalar tuning
cannot fix.

Following their calibration heuristic --- threshold set to ``$1.5\times$ per-step average
regularization cost of a walking gait'' --- we set
\[
  h \;=\; 1.5 \times \bigl(\text{mean per-step regularity cost of the greedy reference solver}\bigr),
\]
i.e.\ calibrated against a known-good reference behaviour rather than in absolute units.

\subsection{A relocation action}

Add a fourth mode, $\MOVE$: relocate the guard nearest to $\xi(a)$ by a bounded
displacement emitted by a second continuous head. Pruning can only remove guards that are
\emph{already} redundant; relocation can \emph{make} a guard redundant by improving a
neighbour. This is the move that greedy fundamentally cannot make and is the main remaining
source of improvement. It is listed as optional because it doubles the continuous action
dimension and should not be introduced until \stageII{} is stable.

\subsection{Hard-instance concentration}

$\Dctx$ reweights to $(0.10, 0.20, 0.35, 0.35)$ over the four tiers, and instances are
additionally sampled in proportion to their recent failure rate --- a simple
prioritised-instance scheme. This is the only place where we use adaptive difficulty; none
of the three reference papers has a performance-gated difficulty curriculum, though it is
standard elsewhere in legged locomotion \cite{rudin2022,narvekar2020curriculum}.

% ===========================================================================
\section{The weight schedule}
% ===========================================================================

\Cref{tab:schedule} is the curriculum. It is deliberately formatted the way
\citet{wang2025dribble} (their Table~I) and \citet{xu2025striker} (their Table~II) format
theirs: one row per reward term, one column per stage.

\begin{table}[t]
\centering
\small
\caption{Reward weights by stage. ``$\to$'' denotes a linear ramp within the stage;
``---'' denotes a term whose triggering action is masked. Terms in the scaffolding group
are the only ones that \emph{must} be scheduled (Document~1, \emph{anneal-iff-non-potential}); every
other schedule entry is an optimisation choice, not a correctness requirement.}
\label{tab:schedule}
\renewcommand{\arraystretch}{1.22}
\begin{tabularx}{\linewidth}{@{}l l X c c c@{}}
\toprule
\textbf{Symbol} & \textbf{Group} & \textbf{Term} & \textbf{\stageI{}} & \textbf{\stageII{}} & \textbf{\stageIII{}} \\
\midrule
$w_{\mathrm{cov}}$   & coverage  & coverage potential (potential-based)        & $1.0$ & $1.0$ & $1.0$ \\
$w_{\mathrm{goal}}$  & coverage  & terminal: $c_T \ge c^{\dagger}$             & $4.0$ & $4.0$ & $4.0$ \\
$w_{\mathrm{fail}}$  & coverage  & terminal: $c_T < c^{\dagger}$               & $2.0$ & $2.0$ & $2.0$ \\
\addlinespace
$w_{\mathrm{par}}$   & parsimony & guard-count potential (potential-based)     & $\mathbf{0.0}$ & $\mathbf{0.02}$ & $0.02$ \\
$w_{\mathrm{prune}}$ & parsimony & $\exp(-\delta^{-}/\sigma_p)$ on $\DEL$      & ---   & $0.15$ & $0.15$ \\
$w_{\mathrm{pars}}$  & parsimony & terminal: $k_{\mathrm{ref}} - k_T$          & $0.0$ & $0.25$ & $0.35$ \\
\addlinespace
$w_{\mathrm{front}}$ & scaffold  & landed in uncovered space                   & $0.05 \to 0$ & $0$ & $0$ \\
$w_{\mathrm{sweep}}$ & scaffold  & boundary-inward ring prior                  & $0.10 \to 0$ & $0$ & $0$ \\
$w_{\mathrm{sep}}$   & scaffold  & minimum guard separation                    & $0.10 \to 0$ & $0$ & $0$ \\
\addlinespace
$w_{\mathrm{proj}}$  & regularity& aimed outside $\Free$                       & $0.25$ & $0.25$ & constraint \\
$w_{\mathrm{void}}$  & regularity& invalid / masked $\DEL$                     & ---    & $0.10$ & constraint \\
$w_{\mathrm{stop}}$  & regularity& $\STOP$ with $c_t < c^{\dagger}$            & ---    & $2.0$  & constraint \\
\midrule
$\beta_{\mathrm{cov}}$ & \multicolumn{2}{@{}l}{advantage weight, coverage critic}  & $1$ & $1$ & $1$ \\
$\beta_{\mathrm{par}}$ & \multicolumn{2}{@{}l}{advantage weight, parsimony critic} & $0$\,\textsuperscript{$\ast$} & $2$ & $2$ \\
$\beta_{\mathrm{reg}}$ & \multicolumn{2}{@{}l}{advantage weight, regularity critic}& $1$ & $1$ & --- \\
\midrule
\multicolumn{3}{@{}l}{Action modes available}                                     & $\ADD$ & $\ADD,\DEL,\STOP$ & $+\,\MOVE$ \\
\multicolumn{3}{@{}l}{Discount $\gamma$}                                          & $0.99$ & $0.997$ & $0.997$ \\
\multicolumn{3}{@{}l}{$\mathcal{L}_{\mathrm{BC}}$ (prune oracle)}                 & --- & $1.0 \to 0$ & $0$ \\
\bottomrule
\end{tabularx}
\\[3pt]
\raggedright\footnotesize
\textsuperscript{$\ast$}\,The parsimony critic is still \emph{trained} in \stageI{} on the
parsimony reward stream; only its advantage weight is zero, so it contributes no policy
gradient. See \cref{sec:shadow}.
\end{table}

Three entries deserve emphasis because they are the ones a reader will want to change first.

\begin{itemize}
  \item \textbf{$w_{\mathrm{par}} = 0.0$ in \stageI{}.} Not $0.005$, not ``small''. Both
        multi-stage reference papers hold their eventual task reward at exactly zero through
        stage 1, and the project brief asks for the same. A small nonzero value is the worst
        of both worlds: too small to shape behaviour, large enough to make ``place no guards''
        locally attractive to a random policy.
  \item \textbf{$w_{\mathrm{cov}} = 1.0$ in every stage.} Because it is potential-based, its
        weight is a pure scaling of an unbiased signal. Nothing is gained by demoting it and
        the risk of forgetting how to cover is real.
  \item \textbf{$\beta_{\mathrm{par}}: 0 \to 2$.} The $2{:}1$ ratio is taken directly from
        \citet{wang2025vision}'s $\Adv_{\mathrm{total}} = 2\Adv_{\mathrm{goal}} + \Adv_{\mathrm{aux}}$.
        The parsimony advantage is small in magnitude but low-variance (the reward is a
        constant $\pm w_{\mathrm{par}}$ plus a bounded bonus), so it tolerates upweighting
        better than the coverage advantage does.
\end{itemize}

% ===========================================================================
\section{Stage transitions}\label{sec:transitions}
% ===========================================================================

\subsection{Performance gates}

Every $E = 10$k steps the current policy is evaluated deterministically on the full held-out
set, per difficulty tier. Let $\mathrm{SR}^{(i)}_\mathrm{easy/mid}$ denote the feasibility
rate ($\Prob[c_T \ge c^{\dagger}]$) at evaluation round $i$, and similarly for
hard/extreme. The gates are:

\[
G_1:\quad
\begin{cases}
  \mathrm{SR}_{\mathrm{easy/mid}}^{(i)} \ge 0.90 & \\
  \mathrm{SR}_{\mathrm{hard/extreme}}^{(i)} \ge 0.75 & \\
  \Ex[c_T] \ge 0.97 &
\end{cases}
\text{for } M = 3 \text{ consecutive } i,
\quad\text{and}\quad w_{\mathrm{front}} = w_{\mathrm{sweep}} = w_{\mathrm{sep}} = 0 .
\]

\[
G_2:\quad
\begin{cases}
  \mathrm{SR}^{(i)} \ge \mathrm{SR}^{\mathrm{end\ of\ \stageI{}}} - 0.02 & \\
  \Ex[k_{\mathrm{peak}} - k_T] \ge 2 & \text{(pruning yield)}\\
  \Ex[k_T / k_{\mathrm{ref}}] \le 1.0 &
\end{cases}
\text{for } M = 3 \text{ consecutive } i .
\]

A hard step cap ($1.5$M for \stageI{}, $1.5$M for \stageII{}) forces the transition if a gate
never fires, with the failure logged. Requiring $M = 3$ consecutive rounds rather than a
single crossing is the difference between a gate and a noise trigger: with only $50$
held-out instances, a success rate near $0.9$ has a binomial standard error of
$\sqrt{0.9 \cdot 0.1 / 50} \approx 4$ percentage points, so a single evaluation carries a
$95\%$ interval more than $16$ points wide. Enlarging the held-out set to $200$ instances
would be a cheaper improvement than tuning $M$.

\subsection{Blending, and the value-function shock}

A hard weight switch changes the reward function, and therefore the return, and therefore
every value target the critic has been fitted to. The critic's error spikes, the advantage
estimates become garbage for several updates, and PPO --- which trusts its advantages ---
takes a large, badly-directed step. This is a plausible contributor to the instability
already logged for this project ($\text{approx\_kl}$ in the tens, $\text{clip\_fraction}
\approx 0.7$--$0.8$, policy collapse after $150$--$200$k steps), and it will get worse, not
better, once a second stage exists. None of the three reference papers discusses it.

Two mitigations, used together:

\begin{enumerate}[label=(\roman*)]
  \item \textbf{Linear blending.} Interpolate every weight linearly over
        $N_{\mathrm{blend}} = 50$k steps:
        $w(t) = (1-\alpha)w^{\mathrm{old}} + \alpha w^{\mathrm{new}}$, $\alpha = t/N_{\mathrm{blend}}$.
        The action mask cannot be blended, so it flips at $\alpha = 0$; the associated
        weights ramp from zero, which means early deletions are nearly free and the policy
        can explore them before they start costing.
  \item \textbf{Tighter trust region across the boundary.} Halve \code{target\_kl} and the
        learning rate for the duration of the blend window, then restore. PPO's KL early-stop
        \cite{schulman2017ppo} is the right tool here and it is already enabled in the current
        code (\code{TARGET\_KL = 0.03}).
\end{enumerate}

\subsection{The shadow parsimony critic}\label{sec:shadow}

Because the reward is grouped (Document~1, multi-critic advantage decomposition), each group has its own
value head, and a head can be trained without being used. During \stageI{} we compute the
parsimony reward stream ($-w_{\mathrm{par}}$ per $\ADD$, plus the terminal parsimony bonus)
with the \emph{\stageII{}} weights, regress $V_{\mathrm{par}}$ against it, and set
$\beta_{\mathrm{par}} = 0$ so it contributes nothing to the policy gradient. At the switch,
$\beta_{\mathrm{par}}$ ramps to $2$ over the blend window against a value head that is already
fitted.

This costs one extra MLP head and no environment changes, and it removes the single largest
source of the transition shock. It has no counterpart in the reference papers, but it is a
direct consequence of adopting \citet{wang2025vision}'s multi-critic decomposition: once the
critics are separate, a critic can be warmed up in advance of the stage that needs it.

% ===========================================================================
\section{Architecture}
% ===========================================================================

\subsection{Networks}

\begin{center}
\small
\begin{tabularx}{\linewidth}{@{}l X@{}}
\toprule
\textbf{Component} & \textbf{Specification} \\
\midrule
Actor trunk   & CNN over the $64\times64\times5$ actor image: $32@8{\times}8$/4, $64@4{\times}4$/2, $64@3{\times}3$/1, ReLU, flatten, FC $256$. Scalar branch: MLP $(64,64)$. Concatenate $\to 320$. \\
Mode head     & Linear $320 \to 3$, categorical over $\{\ADD,\DEL,\STOP\}$ (plus $\MOVE$ in \stageIII{}); invalid modes masked with $-\infty$ logits. \\
Point head    & Linear $320 \to 4$, softplus${}+1$, parameterising two independent $\mathrm{Beta}(\alpha,\beta)$ distributions on $[0,1]$ for $a_x$, $a_y$. \\
Critic trunk  & Separate CNN over the $64\times64\times6$ privileged image plus the full scalar set; MLP $(512,256,128)$ (wider than the actor, as in \citet{wang2025dribble}: $768/256/128$ critic vs.\ $512/256/128$ actor). \\
Critic heads  & Three linear heads $\to 1$: $V_{\mathrm{cov}}$, $V_{\mathrm{par}}$, $V_{\mathrm{reg}}$. \\
Activation    & ELU throughout, following \citet{wang2025vision}. \\
\bottomrule
\end{tabularx}
\end{center}

\subsection{Why a Beta distribution for the point head}

The action is a point in $[0,1]^2$ and the current implementation samples a Gaussian and
clips. This is a poor fit for the AGP specifically, because good guard positions are
frequently \emph{at} the boundary of the action range --- against a wall, in a corner, in the
mouth of a corridor. A clipped Gaussian piles a point mass of probability at the clip
boundary, so the log-probability of a boundary action is wrong, the entropy is wrong, and the
policy gradient at the boundary is biased inward \cite{chou2017beta}. A Beta distribution has
support exactly $[0,1]$, is bias-free at the boundary, and can represent both a sharp
interior mode and a boundary-concentrated one. Given that the \stageI{} sweep prior is
explicitly \emph{encouraging} boundary placements, this is not a minor detail.

\subsection{Symmetry}

The AGP is equivariant under the dihedral group $D_4$ (rotations by multiples of $90^\circ$
and reflections) acting on the rasterised gallery, and under similarity transforms in the
continuous problem. We exploit this twice:

\begin{itemize}
  \item \textbf{Augmentation:} each instance is presented under a uniformly sampled $D_4$
        transform, which multiplies the effective dataset by $8$ at zero cost.
  \item \textbf{Consistency loss:} following the mirror-symmetry loss of
        \citet{wang2025vision} (their Eqs.~5--6, coefficient $10$),
        \[
          \mathcal{L}_{\mathrm{sym}}
          = \Ex\Bigl[\bigl\lVert \xi_\theta(o) - M^{-1}\xi_\theta(M o)\bigr\rVert^2
          + \mathrm{KL}\bigl(\mu_\theta(o)\,\Vert\,\mu_\theta(M o)\bigr)\Bigr],
          \quad M \sim D_4 ,
        \]
        with coefficient $1.0$: the point head should be equivariant, the mode head
        invariant. They needed this because their policy collapsed onto always kicking with
        one leg; the AGP analogue is a policy that collapses onto, say, always sweeping
        clockwise from the top-left.
\end{itemize}

% ===========================================================================
\section{PPO hyperparameters}
% ===========================================================================

\begin{table}[h]
\centering
\small
\caption{PPO settings per stage. The \stageI{} column is close to the current code, with
two changes: vectorised environments and a nonzero entropy coefficient.}
\label{tab:ppo}
\begin{tabularx}{\linewidth}{@{}l c c c X@{}}
\toprule
\textbf{Parameter} & \textbf{\stageI{}} & \textbf{\stageII{}} & \textbf{\stageIII{}} & \textbf{Note} \\
\midrule
Discount $\gamma$        & $0.99$ & $0.997$ & $0.997$ & terminal parsimony must survive $\sim 30$ steps \\
GAE $\lambda$            & $0.95$ & $0.95$  & $0.95$  & \cite{schulman2016gae} \\
Clip range $\epsilon$    & $0.2$  & $0.15$  & $0.10$  & tightened as the policy gets better \\
\code{target\_kl}        & $0.03$ & $0.02$  & $0.015$ & halved during blend windows \\
Learning rate            & $1\!\times\!10^{-4}$ & $5\!\times\!10^{-5}$ & $3\!\times\!10^{-5}$ & constant; halved during blend windows \\
Epochs per update        & $5$    & $5$     & $4$     & as in current code \\
Entropy: mode head       & ---    & $0.01 \to 0.002$ & $0.002$ & critical for $\DEL$ discovery (\cref{sec:oracle}) \\
Entropy: point head      & $0.005$& $0.003$ & $0.001$ & current code uses $0.0$ \\
Parallel envs            & $16$   & $16$    & $16$    & \code{SubprocVecEnv}; currently $1$ \\
Steps per env per update & $256$  & $256$   & $256$   & rollout of $4096$ transitions \\
Minibatches per epoch    & $4$    & $4$     & $4$     & as in \citet{wang2025vision} \\
Max grad norm            & $0.5$  & $0.5$   & $0.5$   & \\
Budget (steps)           & $1.5$M & $1.5$M  & $0.75$M & gates may trigger earlier \\
\bottomrule
\end{tabularx}
\end{table}

The single most impactful line in \cref{tab:ppo} is probably \textbf{parallel envs}. The
current run uses one environment, so a $2048$-step rollout is one long, highly autocorrelated
trajectory through a handful of galleries; the resulting gradient is a poor estimate of the
population gradient, and this is a textbook cause of the enormous \code{approx\_kl} values
($10$--$100$ against a target of $0.01$--$0.05$) already recorded. The reference papers run
$4{,}096$ \cite{wang2025dribble} and $16{,}384$ \cite{wang2025vision} environments; we cannot
match that with Shapely on CPU, but $16$ processes is straightforward and changes the
variance by more than an order of magnitude. Profile the visibility-polygon computation
first --- it is the bottleneck, it is $O(n \log n)$ per guard with $n$ up to $\sim\!100$
vertices including holes, and it is called once per step.

% ===========================================================================
\section{Instance curriculum and augmentation}\label{sec:instances}
% ===========================================================================

The locomotion papers randomise physics (friction, mass, CoM, actuator gains, latency) to
close a sim-to-real gap. We have no sim-to-real gap; our generalisation gap is across
\emph{gallery instances}, so the analogue of domain randomisation is instance randomisation.

\begin{center}
\small
\begin{tabularx}{\linewidth}{@{}l X@{}}
\toprule
\textbf{Mechanism} & \textbf{Detail} \\
\midrule
Difficulty mixture & Tier sampling weights $(0.55,0.30,0.10,0.05) \to$ uniform over \stageI{}; uniform in \stageII{}; $(0.10,0.20,0.35,0.35)$ in \stageIII{}. \\
$D_4$ augmentation & Uniform random rotation/reflection of each instance at reset. \\
Scale jitter       & Uniform scale in $[0.8,1.25]$; all reward terms are already scale-free (areas normalised by $A$, distances by the bounding-box diagonal). \\
Vertex jitter      & Gaussian perturbation of outer vertices with $\sigma = 0.01 L$, rejecting non-simple results; generates fresh instances from a finite dataset. \\
Prioritised sampling & \stageIII{} only: sample instances with probability $\propto (1 + \text{recent failure rate})$. \\
\bottomrule
\end{tabularx}
\end{center}

Vertex jitter deserves a note: with $N_{\mathrm{train}} = 500$ galleries and $1.5$M steps at
$\sim\!10$ steps per episode, each gallery is seen roughly $300$ times, which is well into
memorisation territory. Jitter is the cheapest defence, and it is closer in spirit to what
the locomotion papers do (they never reuse an exact initial state) than a fixed dataset is.

% ===========================================================================
\section{Evaluation and ablations}
% ===========================================================================

\subsection{Baselines}

\begin{itemize}
  \item \textbf{Random interior placement} --- the floor.
  \item \textbf{Greedy max-coverage} over a candidate set (all vertices, reflex vertices, and
        a $32\times32$ interior grid), which defines $k_{\mathrm{ref}}$. This is the baseline
        to beat, and it is a strong one \cite{nemhauser1978,feige1998,ghosh2010}.
  \item \textbf{Greedy $+$ oracle prune} (\cref{sec:oracle}) --- the honest baseline for
        \stageII{}, since our policy has access to the same pruning move.
  \item \textbf{Exact IP} \cite{couto2011,tozoni2016} on the \code{easy} tier only, to
        establish the true optimum where it is computable.
\end{itemize}

Reporting against greedy alone is not enough: if \stageII{} only recovers what a greedy prune
would have given anyway, the learned policy has added nothing.

\subsection{Ablations}

\begin{table}[h]
\centering
\small
\caption{Ablation protocol. Each row isolates one design decision from this document or
Document 1.}
\label{tab:ablations}
\begin{tabularx}{\linewidth}{@{}l X X@{}}
\toprule
\textbf{Ablation} & \textbf{Change} & \textbf{What it tests} \\
\midrule
A1 Single stage      & $w_{\mathrm{par}} = 0.02$ and $\DEL$ unmasked from step $0$ & The core claim. \citet{wang2025dribble} see $86.7\% \to 33.0\%$ from this change. \\
A2 No pruning        & \stageI{} only, run for the full $3$M steps & Isolates \stageII{}'s contribution from extra training. \\
A3 Sparse reward     & Terminal $c_T$ and $k_T$ only, no potentials & Whether the dense telescoping decomposition matters (Document~1, telescoping proposition). \\
A4 No annealing      & Scaffolding weights held constant & Tests the anneal-iff-non-potential rule directly; predicts a wall-hugging, over-separated final policy. \\
A5 Single critic     & One value head on the summed reward & \citet{wang2025vision}'s multi-critic claim. \\
A6 No oracle         & \stageII{} without $\mathcal{L}_{\mathrm{BC}}$ & Whether $\DEL$ is discovered unaided (\cref{sec:oracle}). \\
A7 Hard switch       & $N_{\mathrm{blend}} = 0$, no shadow critic & Measures the transition shock (\cref{sec:shadow}). \\
A8 No sweep prior    & $w_{\mathrm{sweep}} = 0$ throughout & Whether ``edges first, then inwards'' helps or only slows convergence. \\
A9 Gaussian head     & Clipped Gaussian instead of Beta & Boundary-action bias \cite{chou2017beta}. \\
A10 Fixed penalty    & \stageIII{} with fixed coefficients instead of constraints & \citet{xu2025striker}'s inhomogeneous-credit claim. \\
\bottomrule
\end{tabularx}
\end{table}

Metrics are those of Document 1: feasibility rate, mean coverage, guard count conditioned on
success, competitive ratio $k_T/k_{\mathrm{ref}}$, pruning yield, and residual redundancy.
Report all of them per difficulty tier, over at least three seeds --- none of the three
reference papers reports seed variance, which is a weakness worth not inheriting.

% ===========================================================================
\section{Implementation roadmap}
% ===========================================================================

Ordered by dependency, against the existing code.

\begin{enumerate}[label=\textbf{R\arabic*.},leftmargin=2.7em]
  \item \textbf{Fix truncation bootstrapping.} Set \code{info["TimeLimit.truncated"] = True}
        on every \code{truncated} return in \code{ArtGalleryEnv.step}. One line; fixes a
        systematic value-target bias (Document~1, truncation-bootstrapping remark).
  \item \textbf{Add the terminal reward.} $w_{\mathrm{goal}}$, $w_{\mathrm{fail}}$ in
        \code{step()}; $w_{\mathrm{pars}}$ once $k_{\mathrm{ref}}$ exists.
  \item \textbf{Precompute $k_{\mathrm{ref}}$.} Extend \code{generate\_polygon\_dataset.py} to
        run greedy max-coverage per instance and store \code{k\_ref} in the JSON record.
        Also store the greedy guard positions --- they are useful for the oracle and for
        previews.
  \item \textbf{Vectorise.} \code{SubprocVecEnv} with $16$ workers. Profile the ray sweep in
        \code{security.py} first; cache per-guard visibility polygons (already done via
        \code{all\_guards}) and avoid recomputing the union from scratch on $\ADD$.
  \item \textbf{Extend the observation} to the five actor channels and the privileged critic
        channels of Document~1's observation specification. This is the largest change and gates
        everything downstream.
  \item \textbf{Extend the action space} to $\code{Dict}(\code{mode}{:}\,\code{Discrete(3)},
        \code{point}{:}\,\code{Box}(0,1,(2,)))$ with masking. Note that Stable-Baselines3's
        stock PPO does not support action masking or Dict action spaces; either use
        \code{sb3-contrib}'s \code{MaskablePPO} with a factored discrete space, or implement
        a custom \code{ActorCriticPolicy} with a composite distribution. Budget real time
        for this.
  \item \textbf{Weight scheduler.} A callback that owns the stage state machine: current
        stage, blend progress, gate evaluation, and the weight vector pushed to the env via
        \code{set\_attr}. Keep the schedule in one declarative table so it is diffable.
  \item \textbf{Multi-critic.} Three value heads, grouped reward returned in
        \code{info["reward\_groups"]}, advantage combination per Document~1's multi-critic decomposition.
  \item \textbf{Prune oracle and DAgger loss.} \cref{sec:oracle}. Implement the oracle first
        as a standalone evaluation baseline (it is useful on its own), then wire in the BC
        loss.
  \item \textbf{Beta policy head}, $D_4$ augmentation, symmetry loss.
  \item \textbf{\stageIII{}} only after \stageII{} is stable.
\end{enumerate}

Items R1--R4 are worth doing regardless of whether the curriculum is adopted; they are
straightforward fixes to the existing single-stage setup and several of them plausibly
address the instability already observed.

% ===========================================================================
\section{Risks}
% ===========================================================================

\begin{description}[style=nextline,leftmargin=1.2em]
\item[The delete action is never used.] The main risk; \cref{sec:oracle} gives four
  independent mitigations. Diagnostic: log the empirical mode distribution per stage. If
  $\Prob[\DEL] < 0.05$ after $300$k \stageII{} steps, the warm start has failed.
\item[Catastrophic forgetting of placement.] Mitigated by keeping $w_{\mathrm{cov}}$ fixed and
  by the $60\%$ empty-gallery share in \stageII{}. Diagnostic: feasibility rate must not fall
  by more than $2$ points across the boundary --- this is built into gate $G_2$.
\item[Pruning that quietly costs coverage.] The ratchet against the running maximum $\bar c_t$
  prevents drift, but $\epsilon^\star \approx 2.2\%$ (Document 1) still allows a single
  meaningful concession. If evaluation shows coverage creeping down, lower
  $w_{\mathrm{par}}$ rather than raising $w_{\mathrm{prune}}$: the exchange rate
  $\epsilon = w_{\mathrm{par}}/w_{\mathrm{cov}}$ is what sets the tolerance.
\item[Compute.] Every step calls exact ray-sweep visibility, and $\DEL$ additionally requires
  $k$ set-difference evaluations for the redundancy channel. With $16$ envs $\times$ $3$M
  steps this is the dominant cost. Consider computing $\delta^-_j$ on a coarse raster for the
  observation while keeping the exact value for the reward.
\item[Over-fitting to $500$ galleries.] Addressed by vertex jitter and $D_4$ augmentation
  (\cref{sec:instances}); monitor the train/held-out feasibility gap explicitly.
\item[The curriculum is not needed.] Possible: \citet{wang2025vision} solve a comparably
  hard task \emph{single-stage}, and attribute it to potential-based shaping plus a
  multi-critic split --- both of which we also use. Ablation A1 is the test, and it should be
  run early rather than at the end, because a negative result would save most of this work.
\end{description}

% ===========================================================================
\section{Summary}
% ===========================================================================

\begin{enumerate}[label=\textbf{\arabic*.}]
  \item \textbf{\stageI{} (\stageI{})}: cover the gallery with the guard cost at exactly zero,
        under a boundary-inward placement prior and two other scaffolding terms that are all
        annealed to zero before the stage ends. A short over-saturation phase at the end
        deliberately manufactures the redundant guard sets that \stageII{} learns from.
  \item \textbf{\stageII{} (\stageII{})}: unmask $\DEL$ and $\STOP$, switch on the guard cost
        (so adding costs and removing refunds), and add a bounded bonus
        $\exp(-\delta^{-}/\sigma_p)$ that pays in full for a free removal and decays as
        coverage is lost. A coverage ratchet against the running maximum prevents regression.
  \item \textbf{\stageIII{} (\stageIII{}, optional)}: move the regularity penalties into
        state-wise constraints, add a relocation action, and concentrate on hard instances.
  \item The curriculum is entirely a weight schedule (\cref{tab:schedule}) plus an action
        mask over one fixed MDP. Only the three scaffolding terms \emph{must} be annealed;
        everything else is potential-based and unbiased (Document 1).
  \item Three things the reference papers do not supply and we add: quantitative
        stage gates, a blended transition with a shadow-trained parsimony critic, and a
        privileged greedy prune oracle used as a DAgger warm start for the newly unmasked
        delete action.
  \item Run ablation A1 (single-stage) early. It is the one experiment that could invalidate
        the premise.
\end{enumerate}

\bibliographystyle{plainnat}
\renewcommand{\refname}{References}
\bibliography{refs}

\end{document}
